\centerline{\bf \Huge Геометрический разнобой}

\begin{flushright}
        \scriptsize{Эпизод 4. Дождь, после побега}
\end{flushright}

\problem{1}
Пусть $O$ и $H$ - центр описанной окружности и ортоцентр соответственно треугольника $A B C$. Докажите, что одна из площадей $O H A, O H B, O H C$ равна сумме двух других.

\problem{2}
Пусть $A B C$ равнобедренный треугольник, $A C=B C$. Вписанная окружность касается $A B$ в точке $D$ и $B C$ - в точке $E$. Прямая, отличная от $A E$ и проходящая через $A$, пересекает вписанную окружность в точках $F$ и $G$. Прямая $A B$ пересекает $E F$ и $E G$ в точках $K$ и $L$ соответственно. Докажите, что $D K=D L$.

\problem{3}
Пусть $P$ и $Q$ изогонально сопряженные точки в треугольнике $A B C ; P_1 P_2 P_3$ и $Q_1 Q_2 Q_3$ их педальные треугольники. Также пусть $X_1=P_2 Q_3 \cap P_3 Q_2, X_2=P_1 Q_3 \cap P_3 Q_1, X_3= P_1 Q_2 \cap P_2 Q_1$. Докажите, что точки $P, Q, X_1, X_2, X_3$ лежат на одной прямой.

\problem{4}
В треугольнике $A B C$ точки $D, M$ лежат на сторонах $B C$ и $A B$ соответственно, точка $P$ лежит на отрезке $A D$. Прямая $D M$ пересекает $B P, A C$ (вне отрезка), $P C$ (вне отрезка) в точках $E, F$ и $N$ соответственно. Покажите, что если $D E=D F$, то $D M=D N$.

\problem{5}
Пусть $A B C$ остроугольный треугольник, $H$ и $O$ - его ортоцентр и центр описанной окружности соответственно, $M$ и $N$ - середины отрезков $A H$ и $B C$ соответственно. Окружность $\gamma$ с диаметром $A H$ пересекает описанную окружность $A B C$ в точке $G \neq A$, а прямую $A N$ в точке $Q \neq A$. Касательная к $\gamma$ в $G$ пересекает $O M$ в точке $P$. Покажите, что одна из точек пересечения описанных окружностей $\triangle G N Q$ и $\triangle M B C$ лежит на прямой $P N$.

\problem{6}
Пусть $A B C$ остроугольный треугольник с $\angle C=90^{\circ}, H$ - основание высоты из $C$. Точка $D$ выбрана внутри $C B H$ так, что $C H$ делит пополам $A D$. Прямые $B D$ и $C H$ пересекаются в точке $P$. Пусть $\omega$ - полуокружность с диаметром $B D$ пересекающая отрезок $C B . P Q$ - касательная к $\omega$ из точки $P$. Докажите, что прямые $C Q$ и $A D$ пересекаются на $\omega$.